\documentclass[a4paper,11pt]{book}

\usepackage[T1]{fontenc}
\usepackage[utf8]{inputenc}
\usepackage{graphicx}
\usepackage{xcolor}
\usepackage[colorlinks = true, linkcolor = blue]{hyperref}

%quote
\usepackage{csquotes}
\renewcommand{\mkbegdispquote}[2]{\itshape}

%Reset the mathcal
\DeclareMathAlphabet{\mathcal}{OMS}{cmsy}{m}{n}

\usepackage{mathptmx}
\usepackage{amsmath,amssymb,amsthm,textcomp}
\usepackage{multicol}
\usepackage{tikz}
\usepackage[normalem]{ulem}
\usepackage{bbm} % using for the characteristic function

\usepackage{geometry}
\geometry{left=15mm,right=15mm,%
bindingoffset=0mm, top=20mm,bottom=20mm}

\makeatletter
\def\th@plain{%
  \thm@notefont{}% same as heading font
  \itshape % body font
}
\def\th@definition{%
  \thm@notefont{}% same as heading font
  \normalfont % body font
}
\makeatother

\renewcommand{\thechapter}{\Roman{chapter}}
\renewcommand*\thesection{\arabic{section}}

%\newtheorem{theorem}{Theorem}
\newtheorem{theorem}{Theorem}[section]
\renewcommand{\thetheorem}{\arabic{section}.\arabic{theorem}}
\newtheorem{corollary}{Corollary}[theorem]
\newtheorem{lemma}[theorem]{Lemma}
\newtheorem{proposition}[theorem]{Proposition}

\newtheorem{exercise}{Exercise}[subsection]

\theoremstyle{definition}
\newtheorem{definition}{Definition}[section]
\renewcommand{\thedefinition}{\arabic{section}.\arabic{definition}}
\newtheorem*{example}{Example}

\theoremstyle{remark}
\newtheorem*{remark}{Remark}
\newtheorem{review}{Review}

\linespread{1.3}
\newcommand{\linia}{\rule{\linewidth}{0.5pt}}

% my own titles
\makeatletter
\renewcommand{\maketitle}{
\begin{center}
\vspace{2ex}
{\huge \textsc{\@title}}
\vspace{1ex}
\\
\linia\\
\@author \hfill \@date
\vspace{4ex}
\end{center}
}
\makeatother
%%%

%symbol
\AtBeginDocument{
  \let\ran\relax
  \DeclareMathOperator{\dom}{dom}
}\AtBeginDocument{
  \let\ran\relax
  \DeclareMathOperator{\ran}{ran}
}
\DeclareMathOperator{\var}{Var}
\DeclareMathOperator{\cov}{Cov}
\AtBeginDocument{
  \let\div\relax
  \DeclareMathOperator{\div}{div}
}
\DeclareMathOperator*{\esssup}{ess\, sup}
\DeclareMathOperator*{\spa}{span}
\DeclareMathOperator*{\dist}{dist}
\DeclareMathOperator*{\supp}{supp}
\DeclareMathOperator*{\rank}{rank}

\begin{document}
\title{Stochastic -- Finance, Insurance and Population Genetics}

\chapter{Stochastic Finanzmathe}
\section{Discrete Financial Modelling}
\subsection{1-period Financial Modelling}
\subsubsection{Arbitrage and martingale measure}
\begin{definition}[Arbitrage opportunity]
\begin{equation*}
	\exists \bar{\xi} \in \mathbb{R}^{d+1}: \bar{\xi} \bar{\pi} \leq 0, \quad \bar{\xi} \bar{S} \geq 0, a.s., \mathbb{P}(\bar{\xi} \bar{S} > 0) > 0
\end{equation*}
\end{definition}


The next lemma shows we may w.l.o.g. observe arbitrage as portfolio of the risky assets. 
\begin{lemma}
	The following arguments are equivalent 
	\begin{itemize}
		\item there exists an arbitrage $\bar{\xi} \in \mathbb{R}^{d+1}$. 
		\item there exists $\xi \in \mathbb{R}^d: \xi \cdot Y \ge 0$ $\mathbb{P}$ a.s. $\mathbb{P}( \xi \cdot Y > 0) >0$. 
	\end{itemize}
\end{lemma}

\begin{definition}[Martingale measure and equivalent Martingale measure]
A probability measure $\mathbb{P}*$ is martingale if:
\begin{equation*}
	\pi^i = \mathbb{E}^*\left[\frac{S^i}{1 + r}\right] \quad \forall i = 1, \dots, d
\end{equation*}
We denote the set of equivalent martingale measure by:
\begin{equation*}
	\mathcal{P}:= \left\{\mathbb{P}^* \approx \mathbb{P}| \mathbb{E}^*\left[\frac{S^i}{1 + r}\right] \quad \forall i = 1, \dots, d \right\}
\end{equation*}
\end{definition}

\begin{theorem}[The 1st Fundament Theorem of Asset Pricing]
	The market model is free of arbitrage if and only if $\mathcal{P} \neq \emptyset$. In this case there exists a $\mathbb{P}^*$ with $\frac{d\mathbb{P}^*}{d\mathbb{P}} \in L^\infty(\mathbb{P})$
\end{theorem}

We define all the possible payout the the market:
\begin{align*}
	\mathcal{V}:= \{ \bar{\xi}\cdot \bar{S}: \bar{\xi} \in \mathbb{R}^{d +1}\}
\end{align*}
the following theorem suggest if a payout is replicable and the market is free of arbitrage, then it admits only one price.

\begin{theorem}[Law of 1 Price]
	Any replicable payout in a arbitrage-free market has only one price, i.e. given $v \in \mathcal{V}$ with $v = \xi_1 \cdot \bar{S} = \xi_2 \cdot \bar{S}$, $\xi_1, \xi_2 \in \mathbb{R}^{d +1}$, then we have
	\begin{align*}
		\xi_1 \cdot \bar{\pi}  = \xi_2  \cdot \bar{\pi}
	\end{align*}
\end{theorem}
\begin{proof}
	
\end{proof}


\subsubsection{Contingent Claim}
\begin{definition}[Contingent Claim]
A \textbf{claim} is a non-negative RV on the probability space. A \textbf{contingent claim} is a $\sigma(S_1, \dots)$ measurable claim.
\end{definition}

\begin{example}[Put-Call parity]
	C - P = S - K
\end{example}

\begin{theorem}[Arbitrage-free price of a contingent claim]
	Given $\Pi(C)$ the set of all non-negative arbitrage-free price for the contingent claim $C$. If $\mathcal{P} \neq \emptyset$, then:
	\begin{equation*}
		\Pi(C) = \left\{ \mathbb{E}^*\left[\frac{C}{1+r}\right] :  \mathbb{P}^* \in \mathcal{P} \text{ with } \mathbb{E}^*C < \infty \right\}
	\end{equation*}
\end{theorem}
\begin{remark}
The non-negativity of the price of contingent claim is crucial for the non-arbitrage given a martingale measure.\\
\end{remark}

\begin{theorem}[Range of a contingent claim]
	If $\mathcal{P} \neq \emptyset$, and $C$ a contingent claim, then:
	\begin{equation*}
		\pi_{\inf}(C):= \inf \Pi(C) = \inf_{\tilde{P}}\tilde{\mathbb{E}}(\frac{C}{1+r}) = \max \left\{ m \in [0, \infty): \exists \xi \in \mathbb{R}^n:  m +   \xi Y \leq \frac{C}{1 + r} \mathbb{P} f.s. \right\}
	\end{equation*}
	\begin{equation*}
		\pi_{\sup}(C) := \sup \Pi(C) = \sup_{\tilde{P}}\tilde{\mathbb{E}}(\frac{C}{1+r}) = \min \left\{ m \in [0, \infty) : \exists \xi \in \mathbb{R}^n: m + \xi Y \geq \frac{C}{1 + r} \mathbb{P} f.s. \right\} 
	\end{equation*}
\end{theorem}
\begin{remark}\quad\\
\textbf{1.} The fact that the price is a interval doesn't violate the law of 1 price since the contingent claim is in general not replicable. \\
\textbf{2. }The supremum is usually called the minimum super-hedging price. 	
\end{remark}
\begin{proof}
	The proof has the following steps:
	\begin{enumerate}
		\item $\sup_{P^*}\Pi(C) = \inf M$:
		\item $\sup_{P^*}\Pi(C) = \sup_{\tilde{P}}\Pi(C)$: Convex combination of absolute continuous measure and equivalent measure belongs to equivalent measure. We break an absolute continuous measure into these 2 parts therefore it is bounded by the equivalent measure. Then take the $\epsilon$ to 0.  
		\item $\inf M = \min M$:
	\end{enumerate}
\end{proof}

\begin{definition}[Replicable contingent claim]
A contingent claim is replicable, if there is a $\bar{\xi} \in \mathbb{R}^{d+1}$, such that $C = \bar{\xi} \bar{S}$ $\mathbb{P}-a.s.$
\end{definition}
\begin{remark}
By law of one price, we know if that the price is uniquely determined by $\bar{\xi} \bar{S}$	
\end{remark}

\begin{corollary}
	Given $C$ a contingent claim and $\mathcal{P} \neq \emptyset$.
	\begin{enumerate}
		\item $C$ is replicable iff $|\Pi(C)| = 1$
		\item $C$ is not replicable, then $\Pi(C) = (\pi_{\inf}(C), \pi_{\sup}(C))$ is an open interval
	\end{enumerate}
\end{corollary}




\subsubsection{Complete market model.}

\begin{definition}[complete financial market]
A market is complete if every contingent is claim replicable. 
\end{definition}

\begin{theorem}[The 2nd fundamental theorem of Asset pricing]
	A market is complete iff $|\mathcal{P}| = 1$
\end{theorem}
\begin{remark}
If the market is complete, we have:
\begin{equation*}
	V = \{\bar{\xi}\bar{S}|\bar{\xi} \in \mathbb{R}^{d+1}\} = L^0(\Omega, \mathcal{F}, \mathbb{P})\}
\end{equation*}	
\end{remark}

\begin{theorem}
	Every complete market model admits a partition of $\Omega$ in mostly $d+1$ atoms.
\end{theorem}


\subsubsection{1-Period market model with randomised initial price}
\begin{theorem}[The 2nd Fundamental theorem of Asset Pricing, Version II]
The following statements are equivalent
\begin{enumerate}
	\item $K \cap L^0_+ = \{0\}$
	\item $(K \cap L^0_+) \cap L_+^0 = \{0\}$
	\item $\exists \mathbb{P}^* \in \mathcal{P}: \frac{d\mathbb{P}^*}{d\mathbb{P}} \in L^\infty$. 
	\item $\mathcal{P} \neq \emptyset$
\end{enumerate}
\end{theorem}
\begin{proof}
\quad\\
\textbf{2) $\Rightarrow$ 3)}.
\end{proof}


\subsection{Multi-Period Arbitrage theory}
\subsubsection{Arbitrage opportunity and Martingale measures}
\begin{theorem}[The Radon-Nikodym Theorem]
	Given $\mathbb{Q} \ll \mathbb{P}$ and $Z := \frac{d\mathbb{Q}}{d\mathbb{P}}$. Then for all non-negative random variable $X$ and sub-$\sigma$-Algebra $\mathcal{F}_0 \subset \mathcal{F}$:
	\begin{equation*}
		\mathbb{E}^{\mathbb{Q}}(X|\mathcal{F}_0) =\frac{\mathbb{E}^{\mathbb{P}}(XZ|\mathcal{F}_0)}{\mathbb{E}^{\mathbb{P}}(X|\mathcal{F}_0)}  
	\end{equation*}
\end{theorem}

\begin{theorem}[Separating theorem of Hahn-Banach]
\end{theorem}

\begin{definition}[self-financed strategies]
	A trading strategy is called self-financed, if
	\begin{equation*}
		\bar{\xi}_t \bar{S}_t = \bar{\xi}_{t+1} \bar{S}_t \quad \forall t = 1, \dots, T-1
	\end{equation*}
\end{definition}
\begin{remark}
For a trading strategy $\bar{\xi} = (\xi, \xi_0)$, having $\xi$ fixed, we may select $\xi_0$ such that the trading strategy becomes self-financed, in that we choose the $\xi_0$ that satisfies the equation above.  
\end{remark}

\subsubsection*{The Value Process}


\begin{lemma}[Decomposition of the value process]
	The following arguments are equivalent
	\begin{enumerate}
		\item $\bar{\xi}$ is a self-financed strategy
		\item $\bar{\xi}_t \cdot \bar{X_t} = \bar{\xi}_{t + 1} \cdot \bar{X_t}, \forall t$
		\item $V_t = V_0 + \sum_{k = 1}^t \xi_k \cdot (X_k - X_{k - 1})$
	\end{enumerate}
\end{lemma}

\begin{definition}[Martingale measure]
	A martingale measure on $(\Omega, \mathcal{F})$ is a martingale measure, if the discounted price process $X$ is a $\mathbb{Q}$-Martingale.  
\end{definition}

How do we find the equivalent market Martingale measure?

\begin{theorem}
	In a probability space the following arguments are equivalent
	\begin{enumerate}
		\item $\mathbb{P}$ is a martingale measure
		\item For every self-financed strategy $\bar{\xi}$ with $\xi \in L^\infty(\mathbb{P})$, the induced value process is also a martingale
		\item For a self-financed strategy, if the value process $V^-_T$ is integrable(which mean we have bounded capital), then the value process is a martingale
		\item For a self-financed strategy, if the value process is non-negative a.s., then $\mathbb{E}V_T = \mathbb{E}V_0$
	\end{enumerate}
\end{theorem}

\begin{theorem}[The first fundamental theorem of asset pricing]
The multi-period model is arbitrage-free iff $\mathcal{P} \neq \emptyset$.
\end{theorem}
 
\begin{definition}[replicable contingent claim]
\end{definition}

\begin{theorem}
The market is free of arbitrage and $H$ is a discounted contingent claim, then:
\begin{enumerate}
	\item $\Pi(H) = \{ \mathbb{E}^*(H): \mathbb{P}^* \in \mathcal{P} and \mathbb{E}^* < \infty \} \neq \emptyset$
	\item $\pi_{\inf}(H) = \inf_{\mathbb{P}^* \in \mathcal{P}} \mathbb{E}^*(H)$ and $\pi_{\sup}(H) = \sup_{\mathbb{P} \in \mathcal{P}} \mathbb{E}^*(H)$.
\end{enumerate}
\end{theorem}
\begin{proof}
\end{proof}

\begin{definition}[Arbitrage price]
\end{definition}
\begin{remark}
\end{remark}


\subsubsection{Cox-Ross-Rubinstein Modell(Binominal Model)}
\begin{theorem}
	\quad \newline
	The CRR model is arbitrage-free $\Leftrightarrow$ $a < r < b$.
	In this case $\mathcal{P} = \{ \mathbb{P}^*\}$ with
	\begin{equation*}
		p^* := \mathbb{P}^*(R_t = b) = \frac{r - a}{b - a}
	\end{equation*} 
	And under $\mathbb{P}^*$, the returns in each period are independent.
\end{theorem}
\begin{proof}
\end{proof}

\begin{theorem}[Hedging Strategy]
	\begin{align*}
		\xi_t := \Delta_t(S_0, \dots, S_{t - 1}) = \frac{v_t(x_0, x_{t - 1}, x_{t-1}(1 + b)) - v_t(x_0, x_{t - 1}, x_{t-1}(1 + a))}{(1 + r)^{-t} x_{t - 1}(b -a)}
	\end{align*}
	Choose $\xi^0_t$ such that the strategy $\bar{\xi}$ is self-financed with the initial capital $V_0 = \mathbb{E}^*(H)$
\end{theorem}
\begin{proof}
\end{proof}


\subsubsection*{Exotic Derivative}
\begin{definition} \quad
	\begin{enumerate}
		\item running maximum: $M_t : = \max_{0 \leq s \leq t} Z_s$
	\end{enumerate}
\end{definition}

\begin{theorem}[Reflexion Principle]
	$Z_t$ standard random walk. $\mathbb{P}$ uniform distribution.
	$\forall k \in \mathbb{N}, l \in \mathbb{N}_0$,
	\begin{equation*}
		\mathbb{P}(M_T \geq k, Z_T = k - l) = \mathbb{P}(Z_t = k + l)
	\end{equation*}
\end{theorem}
\begin{proof}
	By using the obvious fact
	\begin{equation*}
		k+l > k, \quad l \geq 0
	\end{equation*}
	to get rid of $M_T \geq k$. 
\end{proof}

\begin{theorem}[Reflexion Principle under $\mathbb{P}^*$]
\begin{equation*}
	\mathbb{P}^*[M_T \geq k, Z_T = k - l] = \left( \frac{1-p^*}{p^*} \right)^l \mathbb{P}^*[Z_T = k + l]
\end{equation*}
\end{theorem}
\begin{proof}
	Calculate the bounded density: $\frac{d\mathbb{P}^*}{d\mathbb{P}}$
	\begin{equation*}
		\mathbb{P}^*(\omega) = (p^*)^{[T + Z_T(\omega)]/2}(1-p^*)^{[T - Z_T(\omega)]/2} = \frac{d\mathbb{P}^*}{d\mathbb{P}} d\mathbb{P}(\omega) = \frac{d\mathbb{P}^*}{d\mathbb{P}} 2^{-T} 
	\end{equation*}
\end{proof}

\begin{theorem}[Up-and-in call option]
\end{theorem}
\begin{proof} \quad
	\begin{enumerate}
		\item Split the expectation
		\item reflection principle
		\item measure change
	\end{enumerate}
\end{proof}

\subsubsection*{Black-Schole Convergence}
\begin{lemma}[Central limit theory for sum]
\end{lemma}

\begin{lemma}[Slusky lemma]
\end{lemma}

\begin{theorem}[Black-Schole Convergence theorem]
Assuming
\begin{enumerate}
	\item the returns are i.i.d., (clear from 2.2.1)
	\item $\sum_{k = 1}^N Var^*(R_N^k) \rightarrow T\sigma^2$
	\item $-1 \leq \alpha_N \leq R^N_t \leq \beta_N$ with $\lim \alpha_N = \lim \beta_N = 0$
\end{enumerate}
We have:
\begin{equation*}
	S^{(N)}_{N} \overset{d}{\longrightarrow} S_T:= S_0\exp\left((rT - \frac{1}{2}\sigma T) + \sigma W_T \right) \quad W_t \sim N(0, T) 
\end{equation*}
\end{theorem}

\begin{exercise}
	down-and-in put option
\end{exercise}

\begin{exercise}
	\begin{equation*}
		C_{BS}^{(N)}
	\end{equation*}
	Determine the price in the $N^{th}$ CRR-Model at $t = 0$. Calculate the limit as $N \rightarrow \infty$. 
\end{exercise}


\subsubsection{American option}




\newpage
\section{Time-continuous finance modelling}
\subsection{Ito-calculus}
\begin{definition}[Brownian motion]
	A stochastic process $(W_t)_{t \\geq 0}$, 
	\begin{enumerate}
		\item $W_0 = 0$ $\mathbb{P}$-$a.s.$
		\item independent increment:
		\item normal distributed increment:
		\item $\mathbb{P}$-a.s. continuous path
	\end{enumerate}
\end{definition}
\begin{remark}
\begin{equation*}
	\cov(W_s, W_t) =  
\end{equation*}	
\end{remark}

\begin{theorem}
	 The following are martingales:
	 \begin{enumerate}
	 	\item $(W_t)$
	 	\item $(W_t^2 - t)$
	 	\item $\left(\exp(\sigma W_t - \frac{\sigma^2}{2}t)\right)$
	 \end{enumerate}
\end{theorem}
\begin{proof}\quad
	\begin{itemize}
		\item The second
		\begin{align*}
			\mathbb{E}[W_t^2 - t|\mathcal{F}_s] &= \mathbb{E}[(W_t - W_s)^2 - W_s^2 - 2W_sW_t|\mathcal{F}_s] - t\\
			&= \mathbb{E}(W_t-W_s)^2 + W_s^2 - t\\
			&= W_s^2 - s
		\end{align*}
		 
		\item The third
		\begin{align*}
			\mathbb{E}[e^{\sigma W_t - \frac{\sigma^2}{2}t} |\mathcal{F}_s] &= e^{\sigma W_s - \frac{\sigma^2}{2}t}\mathbb{E}[e^{\sigma(W_t - W_s)}]\\
			&= e^{\sigma W_s - \frac{\sigma^2}{2}t}\mathbb{E}[e^{x\sigma \sqrt{t-s}}] \quad \text{($x \sim N(0, 1)$)}\\
			&= e^{\sigma W_s - \frac{\sigma^2}{2}s}
		\end{align*}
	\end{itemize}
\end{proof}


\begin{definition}[total variation]
	\begin{equation*}
		TV_t(X) := \sup_{\prod}\sum_{t_i \leq t}|X_{t_{i+1} \wedge t} - X_{t_1}]
	\end{equation*}
	a continuous and deterministic path $X: [0, \infty) \rightarrow \mathbb{R}$ is of 
	\begin{enumerate}
		\item  finite variation: $\forall t \in [0, \infty): TV_t < \infty$.
		\item bounded variation for $T \in [0, \infty)$: $TV_T < \infty$. 
	\end{enumerate}
\end{definition}

\begin{lemma}
	$X$ is of finite variation $\Leftrightarrow$ $\exists A^+, A^- \text{ monoton}: X = A^+ - A^-$
\end{lemma}
\begin{remark}
Lebeague theory and define a integral $\int_0^t f(s) dX_s $	
\end{remark}

\begin{definition}[Quadratic variation]
	$X_t: [0, \infty) \rightarrow \mathbb{R}$ continuous. For $\Pi_n \rightarrow 0, t \rightarrow \langle X \rangle_t$ continuous, we define the quadric variation of $(X_t)$ as:
	\begin{equation*}
		\langle X \rangle_t := \lim_n \sum_{t_i < t, t_i \in \Pi_n} (X_{t_{i+1} - t_i} - X_{t_1})^2 
	\end{equation*}
\end{definition}
\begin{remark}
$\langle X \rangle_t$ depends on $\Pi_n$.	
\end{remark}

\begin{lemma}
	$X$ is of bounded variation for $T$ $\Rightarrow$ $\langle X \rangle_T = 0$ 
\end{lemma}

\begin{corollary}
	$A$ is of bounded variation $\Rightarrow$ $\langle X + A \rangle_T = \langle X \rangle_T$
\end{corollary}

\begin{lemma}
$\langle X \rangle_t(\Pi_n)$	 is well defined, then $\forall g \in C(\mathbb{R}_+, \mathbb{R})$
\begin{equation*}
	\sum_{t_i \in \Pi_n t_i < t} g(t_i) (X_{t_{i+1}} - X_{t_i})^2 \rightarrow \int_0^t g(s)d\langle X \rangle_s
\end{equation*}
\end{lemma}
\begin{proof}
	$\langle X \rangle_t$ monotone and of finite variation($\sigma$-finite measure)\\
	$\Rightarrow$ $\mu^n(0, t) := \sum_{t_i \leq t, t_i \in \Pi_n} (X_{t_{i+1}\wedge t} - X_{t_i})^2$  defines a measure. And since the quadratic variation is well-defined it converges pointwise to the measure induced by $\langle X \rangle$.\\
	$\Rightarrow$ It implies the weak convergence and by the definition of weak convergence we have the result.
\end{proof}

\begin{lemma}[Excersice]
	$X$ of finite quadratic variation, then
	\begin{equation*}
		\forall g \in C^1 \forall t \in [0, T]: \langle g(X) \rangle_t := \sum_{t_i \leq t, t_i \in \Pi_n} (g(X_{t_{i+1}\wedge t}) - g(X_{t_i}))^2 < \infty
	\end{equation*} 
	and 
	\begin{equation*}
		\langle g(X) \rangle_t = \int_0^t g'(X_s)^2d\langle X \rangle_s
	\end{equation*}
\end{lemma}
\begin{proof}
By mean value theorem, for $\forall 0< s < t < T \exists o \in (s, t):$
\begin{align*}
	g(X_t) - g(X_s) &= g'(X_s)(X_t - X_s) + (g'(X_o) - g'(X_s))(X_t - X_s)\\
	(g(X_t) - g(X_s))^2 &= g'(X_s)^2(X_t - X_s)^2 + [2g'(X_s) (g'(X_o) - g'(X_s)) + (g'(X_o) - g'(X_s))^2](X_t - X_s)^2\\
	&= g'(X_s)^2(X_t - X_s)^2 + [(g'(X_s) - g'(X_o))(g'(X_o) - g'(X_s))^2](X_t - X_s)^2\\
	&\Rightarrow g'(X_s)^2(X_t - X_s)^2 \quad \text{by the uniform continuity of g'(x) on the interval $[s, t]$}
\end{align*}
\end{proof}

\begin{theorem}
	$(W_t)_{0, T}$ Brownian motion, $\Pi_n \rightarrow 0$ partition of $[0, T]$ and $\sum |\Pi_n| < \ infty$, then
	\begin{equation*}
		\lim_n \sum (W_{t_{i+1}} - W_{t_i})^2 \rightarrow t \quad \forall t \in [0, T] \quad \mathbb{P}\text{-a.s.} 
	\end{equation*}
\end{theorem}
\begin{proof}
	
\end{proof}

\begin{theorem}[Ito-Formal]
$X: [0, T] \rightarrow \mathbb{R}$ continuous and admits continuous quadratic variation. $A:[0, T] \rightarrow \mathbb{R}$ continuous and admits bounded total variation. Then $\forall F\in C^{1, 2}(\mathbb{R}\times \mathbb{R}, \mathbb{R})$
\begin{equation*}
	\int_0^t F_x(A_s, X_d) dX_s := \lim \sum F(A_{t_i}, X_{t_i})(X_{t_{i+1} \wedge t_i} - X_{t_i})
\end{equation*}
is well defined and
\begin{equation*}
	F(A_t, X_t) = F(X_0, A_0) + \int_0^t F_a(A_s, X_s) dA + \int_0^t F_x (A_s, X_s) dX_s + \int_0^t F_{xx} (A_s, X_s) d\langle X \rangle_s
\end{equation*}
\end{theorem}

\begin{corollary}
\end{corollary}

\begin{lemma}[Associativity of Ito-formel]
	$X: \mathbb{R} \rightarrow \mathbb{R} \in C(\mathbb{R})$ and $\forall t: \langle X \rangle_t < \infty $, $f \in C^2(\mathbb{R}), Y: \mathbb{R}_+ \rightarrow \mathbb{R}$ continuous, Assuming the integral exists:
	\begin{equation*}
		\int_0^t Y_s f'(X_s)dX_s = \lim_n \sum_{t_i < t, t_i \in \Pi_n} Y_{t_i}f'(X_{t_i}) (X_{t_{i+1} - t_i} - X_{t_1})
	\end{equation*}
	and 
\end{lemma}


\begin{example}
\end{example}

\begin{example}	
\end{example}

\begin{example}[The geometric Brownian motion]
\begin{equation*}
	S_t = s_0 e^{\sigma W_t + (\mu - \frac{\sigma^2}{2})t}
\end{equation*}
\end{example}


\subsection{Perfect dynamic hedge with out interest}
\begin{example}[Black-Schole Model]	
\end{example}

\begin{example}[Bachelier Model] \quad
\begin{enumerate}
	\item $X_t = X_0 + \mu t + \sigma W_t$
	\item $\sigma(s, x) = \sigma$ constant
\end{enumerate}
\end{example}

\subsubsection*{Assumption(Volatility profile):}
\begin{itemize}
	\item The quadratic variation of the price $X_t$ could be expressed as
\begin{equation*}
	\langle X \rangle_t := \int_0^t \sigma(s, X_s) ds
\end{equation*}
	\item 
\end{itemize}


\subsubsection*{The Heat equation}
\begin{lemma}
	The boundary value problem:
	\begin{align}
		F_t(t, x) + \frac{1}{2} \sigma^2(x, t) F_{xx}(t, x) &= 0\\
		F(T, x) &= f(x)
	\end{align}
	If $F \in C^{1, 2}((0, T) \times \mathbb{R}) \cap C([0, T] \times \mathbb{R})$ and the BVP admits a solution then the value process
	\begin{equation*}
		H(\omega) := f(X_T(\omega)) = F(0, X_0) + \int_0^T F_x (s, X_s(\omega)) dX_s(\omega)
	\end{equation*}
	exist $\mathbb{P}$-a.s.
\end{lemma}

\begin{theorem}
	$f$ continuous and satisfies
	\begin{equation*}
		|f(x)| \leq \alpha ( 1 + e^{k|x|^\beta}) \quad \alpha , \beta < 0, k \in (0, 2)
	\end{equation*}
	The the BVP 
	\begin{align}
		F_t(t, x) + \frac{1}{2} \sigma^ F_{xx}(t, x) &= 0, x \in [0, T)\\
		F(T, x) &= \lim_{t \rightarrow T}f(x)
	\end{align}
	has the solution
	\begin{equation*}
		F(t,x) := \int_\mathbb{R} f(y) P(\sigma^2(T - t), x-y)dy
	\end{equation*}
\end{theorem}



\chapter{Interest Rate Modelling}
\section{Introduction}
\section{Stochastic basics}

\begin{theorem}[Bayes's rule for conditional expectation]
Let 
\[
R = \frac{d\mathbb{\hat{P}}}{d\mathbb{P}}
\]
be the Radon–Nikodym derivative associated with 
\((\Omega, \mathcal{F}, \mathbb{P})\) and \((\Omega, \mathcal{F}, \mathbb{\hat{P}})\),
and \(X\) a random variable. Then
\[
\mathbb{E}^{\mathbb{\hat{P}}}[X \mid \mathcal{F}_t] = 
\frac{\mathbb{E}^{\mathbb{P}}[R X \mid \mathcal{F}_t]}{\mathbb{E}^{\mathbb{P}}[R \mid \mathcal{F}_t]}.
\]
\end{theorem}

\begin{theorem}[It\^o's Isometry]
\label{thm:ito-isometry}
Let $\{W_t\}_{t \ge 0}$ be a standard Brownian motion on a probability space 
$(\Omega, \mathcal{F}, \{\mathcal{F}_t\}, \mathbb{P})$. 
Suppose $X = \{X_s\}_{s \in [0,T]}$ is an \emph{adapted}, 
square-integrable process (predictable in the strict sense) such that
\[
   \mathbb{E}\!\Bigl[\int_0^T X_s^2 \, ds\Bigr] < \infty.
\]
Then the It\^o integral
\[
   I_T(X) \;=\; \int_0^T X_s \, dW_s
\]
satisfies
\[
   \mathbb{E}\Bigl[\bigl(I_T(X)\bigr)^2\Bigr]
   \;=\;
   \mathbb{E}\Bigl[\int_0^T X_s^2 \, ds\Bigr].
\]
\end{theorem}

\begin{proof}[Sketch of Proof]
We prove the result in two main steps:

\paragraph{Step 1: The case of simple (step) processes.}
Let $X_t$ be a simple, predictable process of the form
\[
  X_t(\omega) \;=\; \sum_{k=0}^{n-1} \xi_k(\omega)\,\mathbf{1}_{[t_k, t_{k+1})}(t),
\]
where $0 = t_0 < t_1 < \cdots < t_n = T$ and each $\xi_k$ is 
$\mathcal{F}_{t_k}$-measurable. Then the It\^o integral is
\[
  \int_0^T X_s \, dW_s 
  \;=\; \sum_{k=0}^{n-1} \xi_k \bigl(W_{t_{k+1}} - W_{t_k}\bigr).
\]

Since $W_{t_{k+1}} - W_{t_k}$ are independent increments with mean zero 
and variance $t_{k+1}-t_k$, we have
\[
  \mathbb{E}\Bigl[\Bigl(\sum_{k=0}^{n-1} 
    \xi_k \bigl(W_{t_{k+1}} - W_{t_k}\bigr)\Bigr)^2\Bigr]
  \;=\;
  \sum_{k=0}^{n-1} 
    \mathbb{E}\bigl[\xi_k^2\bigr]\;\mathbb{E}\bigl[(W_{t_{k+1}} - W_{t_k})^2\bigr].
\]
The cross terms vanish because of independence and zero mean of each increment. 
Using $\mathbb{E}[(W_{t_{k+1}} - W_{t_k})^2] = t_{k+1} - t_k$,
\[
  \sum_{k=0}^{n-1} \mathbb{E}[\xi_k^2]\,(t_{k+1} - t_k)
  \;=\;
  \mathbb{E}\Bigl[ \sum_{k=0}^{n-1} \xi_k^2\,(t_{k+1} - t_k) \Bigr]
  \;=\;
  \mathbb{E}\Bigl[\int_0^T X_s^2 \, ds\Bigr].
\]
Hence, for simple processes $X_t$,
\[
  \mathbb{E}\Bigl[\Bigl(\int_0^T X_s \, dW_s\Bigr)^2\Bigr]
  \;=\;
  \mathbb{E}\Bigl[\int_0^T X_s^2 \, ds\Bigr].
\]

\paragraph{Step 2: Approximation by simple processes.}
For a general (predictable) process $X_t$ satisfying 
$\mathbb{E}[\int_0^T X_s^2\,ds] < \infty$, we can approximate $X$ in $L^2$ 
by a sequence of simple processes $\{X^{(n)}\}$ such that
\[
   \int_0^T \mathbb{E}\bigl[(X_s - X_s^{(n)})^2\bigr]\,ds 
   \;\xrightarrow[n\to\infty]{}\; 0.
\]
Define 
\[
   I_T(X^{(n)}) \;=\; \int_0^T X_s^{(n)} \, dW_s.
\]
By the result for simple processes (Step~1),
\[
   \mathbb{E}\Bigl[\bigl(I_T(X^{(n)})\bigr)^2\Bigr]
   \;=\;
   \mathbb{E}\Bigl[\int_0^T \bigl(X_s^{(n)}\bigr)^2 \, ds\Bigr].
\]
One shows, using It\^o isometry itself on the difference and 
dominated convergence arguments, that
\[
   \int_0^T X_s^{(n)} \, dW_s 
   \;\xrightarrow{L^2}\;
   \int_0^T X_s \, dW_s
\]
and
\[
   \int_0^T \bigl(X_s^{(n)}\bigr)^2 \, ds
   \;\xrightarrow{L^1}\;
   \int_0^T X_s^2 \, ds.
\]
Thus, taking limits yields
\[
   \mathbb{E}\Bigl[\Bigl(\int_0^T X_s \, dW_s\Bigr)^2\Bigr]
   \;=\;
   \mathbb{E}\Bigl[\int_0^T X_s^2 \, ds\Bigr],
\]
as required.
\end{proof}

\begin{theorem}[Ito's Lemma in One Dimension]
\label{thm:Ito}
Let $X_t$ be a one-dimensional Ito process satisfying the stochastic differential equation
\[
   dX_t = \mu(t, X_t)\,dt \;+\; \sigma(t, X_t)\, dW_t,
\]
where $W_t$ is a standard Brownian motion. Suppose $f(t, x)$ is a function such that
\[
   f \in C^{1,2}([0,T] \times \mathbb{R}) 
   \quad(\text{once in } t, \text{twice in } x).
\]
Then $Y_t = f(t, X_t)$ satisfies
\[
   df(t, X_t)
   \;=\;
   \left(
     \frac{\partial f}{\partial t} 
     + \mu \,\frac{\partial f}{\partial x} 
     + \tfrac{1}{2}\,\sigma^2 \,\frac{\partial^2 f}{\partial x^2}
   \right)\!dt
   \;+\;
   \sigma \,\frac{\partial f}{\partial x}\,dW_t.
\]
\end{theorem}

\begin{proof}
\textbf{Step 1: Setup and Taylor expansion.}

Fix a small time increment $\Delta t > 0$. Let
\[
   \Delta X_t \;=\; X_{t+\Delta t} - X_t 
   \;\approx\; \mu(t, X_t)\,\Delta t \;+\; \sigma(t, X_t)\,\Delta W_t,
\]
where $\Delta W_t = W_{t+\Delta t} - W_t \sim \mathcal{N}(0,\Delta t)$.

Consider 
\[
   \Delta f \;=\; f(t + \Delta t, X_t + \Delta X_t) \;-\; f(t, X_t).
\]
By a second-order Taylor expansion around $(t, X_t)$,
\[
   \Delta f 
   \;\approx\;
   \frac{\partial f}{\partial t}(t,X_t)\,\Delta t
   \;+\;
   \frac{\partial f}{\partial x}(t,X_t)\,\Delta X_t
   \;+\;
   \tfrac{1}{2}\,\frac{\partial^2 f}{\partial x^2}(t,X_t)\,(\Delta X_t)^2.
\]
Higher-order (small) terms are omitted as $\Delta t \to 0$.

\medskip
\textbf{Step 2: Substitute the SDE and use $(dW_t)^2 \approx dt$.}

Since
\[
   \Delta X_t 
   \;\approx\;
   \mu\,\Delta t + \sigma\,\Delta W_t,
\]
we have
\[
   (\Delta X_t)^2
   \;=\;
   (\mu\,\Delta t + \sigma\,\Delta W_t)^2
   \;\approx\;
   \mu^2 (\Delta t)^2 + 2\,\mu\,\sigma\,\Delta t\,\Delta W_t + \sigma^2\,(\Delta W_t)^2.
\]
As $\Delta t \to 0$:
\[
   (\Delta t)^2 \;\to\; 0 
   \quad\text{(faster than } \Delta t\text{)},
   \quad
   \Delta t \,\Delta W_t \;\text{is } O((\Delta t)^{3/2}),
   \quad
   (\Delta W_t)^2 \;\approx\; \Delta t.
\]
Hence the dominant term in $(\Delta X_t)^2$ is $\sigma^2 \Delta t$.

Substitute back:
\[
   (\Delta X_t)^2 \;\approx\; \sigma^2 (\Delta W_t)^2 \;\approx\; \sigma^2 \,\Delta t.
\]

\medskip
\textbf{Step 3: Collect terms in $\Delta t$ and $\Delta W_t$.}

Therefore,
\[
   \Delta f 
   \;\approx\;
   \frac{\partial f}{\partial t}\,\Delta t
   \;+\;
   \frac{\partial f}{\partial x}\,\bigl(\mu\,\Delta t + \sigma\,\Delta W_t\bigr)
   \;+\;
   \tfrac{1}{2}\,\frac{\partial^2 f}{\partial x^2}\,\sigma^2 \Delta t.
\]
Combine the terms:
\[
   \Delta f
   \;=\;
   \Bigl(\frac{\partial f}{\partial t} 
         + \mu\,\frac{\partial f}{\partial x} 
         + \tfrac{1}{2}\,\sigma^2\,\frac{\partial^2 f}{\partial x^2}\Bigr)\,\Delta t
   \;+\;
   \sigma\,\frac{\partial f}{\partial x}\,\Delta W_t.
\]

\medskip
\textbf{Step 4: Pass to the limit as $\Delta t \to 0$.}

In the limit $\Delta t \to 0$, we write in differential form:
\[
   df(t,X_t)
   \;=\;
   \Bigl(\frac{\partial f}{\partial t} 
         + \mu\,\frac{\partial f}{\partial x} 
         + \tfrac{1}{2}\,\sigma^2\,\frac{\partial^2 f}{\partial x^2}\Bigr)\,dt
   \;+\;
   \sigma\,\frac{\partial f}{\partial x}\,dW_t.
\]

This is precisely the claimed \emph{Ito's Lemma}, which generalizes the usual chain rule to the stochastic setting, accounting for the quadratic variation of Brownian motion.
\end{proof}

\begin{remark}
The extra term $\tfrac{1}{2}\,\sigma^2\,\frac{\partial^2 f}{\partial x^2}$ in the $dt$ component
appears because $(dW_t)^2 \approx dt$. This sets Ito's Lemma apart from the standard
deterministic chain rule.
\end{remark}


\subsection{Financial Market Basics}
\begin{definition}[Financial Market]
We assume \( p \) (dividend-free\(^1\)) assets \( X(t) = [X_1(t), \dots, X_p(t)]^\top \), which are driven by Ito processes:
\[
dX(t) = \mu(t, \omega)\,dt + \sigma(t, \omega)\,dW(t).
\]
\end{definition}

\begin{definition}[Trading Strategy]
A trading strategy represents a \textcolor{orange}{predictable} adapted process (of asset weights)
\[
\phi(t, \omega) = [\phi_1(t, \omega), \dots, \phi_p(t, \omega)]^\top.
\]
The value of the trading strategy (or corresponding portfolio) is
\[
\pi(t) = \phi(t)^\top X(t).
\]
\end{definition}


\begin{definition}[Numeraire]
	A Numeraire is a positive asset $N(t)$ of our market. An equivalent martingale measure(w.r.t. numeraire $N(t)$) is a measure $\mathbb{Q}$ such that the normalised asset price $\frac{X_i(t)}{N(t)}$ is a $\mathbb{Q}$ martingale.
	The Radon-Nikodym density from $\mathbb{Q}^N$ to $\mathbb{Q}^M$ is
\[
\zeta(t) = \frac{M(t)}{N(t)} \cdot \frac{N(0)}{M(0)}.
\]
	
\end{definition}

\begin{theorem}[The Fundamental Theorem of Asset Pricing]
	
\end{theorem}



\section{Basic Fixed Income Modelling}
In the next we introduce a central concept in the financial market. 
\begin{definition}[Short rate and (abstract)bank account]
	An asset with price $B(t)$ given by $B(0) = 1$ and 
	\begin{equation*}
		dB(t) = r(t)B(t)dt
	\end{equation*}
	therefore
	\begin{equation*}
		B(t) = \exp\left(\int_0^t r(s) ds\right)
	\end{equation*}
\end{definition}

The short rate is normally considered as \textbf{risk-free rate}, set by the central bank. It is crucial to determine the equivalent Martingale measure. Accordingly, the most relevant assets are Zero Coupon Bonds, the short rate determine the most basis asset in financial market: money asset account, which serves as numeraire of financial assets).


\begin{definition}[(ZCB)zero coupon bond]
Define 
\begin{equation}
	P(t, T) 
\end{equation}
the price of a zero coupon bond at time $t$ of maturity $T$ with the payout at terminal time $1$($P(T, T) = 1$). It is a martingale with respect to w.r.t. $B(t)$ as Numeraire. 
	\begin{equation*}
		\mathbb{E}^{\mathbb{Q}}[ \frac{P(t, T)}{B(t)}] = \mathbb{E}^\mathbb{Q}\left[\frac{P(T, T)}{B(T)}\right] = \mathbb{E}^\mathbb{Q} \left[\exp(-\int_0^T r(s)ds)\right].
	\end{equation*} 
	Multiplying both sides by $B(t)$, it holds:
	\begin{equation*}
		\mathbb{E}^\mathbb{Q} [P(t, T)] = \mathbb{E}^\mathbb{Q}\left[\exp(-\int_t^T r(s) ds)\right]
	\end{equation*}
\end{definition}
 

\begin{definition}[Discounted Cashflow Methodology]
Given $V$ the cash flow stream, we have by FTAP
	\begin{equation*}
		\frac{V(t)}{B(t)} = \sum_{i = 1}^N \mathbb{E}^{\mathbb{Q}}\left[\frac{V_i}{B(T_i)}|\mathcal{F}_t\right] 
	\end{equation*}
	We denote 
	\[
	\mathbb{E}^{T_i}[ \cdot ]
	\]
	the expectation in $T_i$ forward measure with zero coupon bond numeraire $P(t, T_i)$. Do the measure change
	\[
\frac{V(t)}{B(t)} = \sum_{i=1}^{N} \mathbb{E}^{T_i} \left[ \frac{P(t, T_i)}{B(t)} \cdot \frac{V_i}{P(T_i, T_i)} \;\middle|\; \mathcal{F}_t \right].
\]
We have therefore
	\begin{equation*}
		V(t) = \sum^N_{i = 1}P(t, T_i) \mathbb{E}^{T_i}[V_i|\mathcal{F}_t]
	\end{equation*}
	If the future cash flow is known, then
	\begin{equation*}
		V(t) = \sum^N_{i = 1}P(t, T_i) V_i
	\end{equation*}
	In general, the challenge lies in predicting the future cash flow.
\end{definition}



\section{Yield Curve and Linear Product}
\subsection{Static Yield Curve and Market Convention}

\begin{definition}[Yield Curve]	
A yield curve at an observation time $t$ is the function of the \textbf{zero coupon bonds} $P(t, \cdot): [t, \infty) \rightarrow \mathbb{R}^+$ for maturities $T \geq t$ to the IRR. The x-axis is the maturity, the y-axis is the calculated interest rate $z$. There are different methods to calculate $z$:  \\
\textbf{discrete compounded zero rate curve}
\begin{equation*}
	P(t, T) = ( 1 + \frac{z_p(t, T)}{p})^{-p(T-t)}
\end{equation*}
\textbf{simple compounded zero rate curve}
\begin{equation*}
	P(t, T) = \frac{1}{1 + z_0(t, T)\cdot(T - t)}
\end{equation*}
\textbf{continuously compounded zero rate curve}
\begin{equation*}
	P(t, T) = \exp\{ - z(t, T) \cdot (T - t) \}
\end{equation*}
\end{definition}
\begin{remark}
Typically, we have $t = 0$. 	
\end{remark}


\begin{definition}[Continuous Forward Rate]
\begin{equation*}
	f(t, T) = -\frac{\partial ln(P(t, T))}{\partial T}
\end{equation*}
\end{definition}
\begin{remark}
\quad
\begin{enumerate}
	\item The continuous rate is the spot rate at the maturity. 
	\item Holiday Calendar. Business day Convention.
\end{enumerate}
\end{remark}

\begin{definition}[Day Count Convention]
A day count convention is a function 
\[
\tau : \mathcal{D} \times \mathcal{D} \to \mathbb{R}
\]
which measures a time period between dates in terms of years.
\end{definition}


\subsection{Multi-Curve Discounted Cash Flow Pricing}
By simple calculation we determine the fair rate of the interbank lending rate:
\begin{definition}[Spot Libor Rate]
Fair rate for interbank lending deal with trade date $T$.
\begin{equation*}
	L(T;T_0, T_1) = \bigg[\frac{P(T, T_0)}{P(T, T_1)} - 1\bigg] \frac{1}{\tau}
\end{equation*}
\end{definition}
\begin{remark}
\quad
\begin{enumerate}
	\item The Libor rate is daily updated.
	\item The interval $[T_0, T_1]$ is typically 1m, 3m, 6m, 12m. 
\end{enumerate}	
\end{remark}

We observe the plain vanilla legs with DCF methodology: 
\[
V(t) = \sum_{i=1}^{N} P(t, T_i) \cdot \mathbb{E}^{T_i} \left[ L(T_{i-1}^F, T_{i-1}, T_i) \cdot \tau_i \mid \mathcal{F}_t \right]
\]



The next theorem shows: Libor is a martingale in the terminal measure with the Numeraire of zero coupon bond of maturity $T_1$.

\begin{theorem}[Martingale Property of Libor rate]
	The Libor rate $L(T, T_0, T_1)$ is a martingale with respect in the $T_1$-forward measure. And we have:
	\begin{equation*}
		\mathbb{E}^{T_1} [L(T;T_0, T_1)|\mathcal{F}_t] = \left[\frac{P(t, T_0)}{P(t, T_1)} - 1\right] \frac{1}{\tau} = L(t; T_0, T_1)
	\end{equation*}
	Which allows us to price the Libor leg based on $t$:
	\begin{equation*}
		V(t) = \sum^N_{i = 1} P(t, T_i) 
	\end{equation*}
\end{theorem}
\begin{proof}
Fair Libor rate at fixing time $T$ is
\[
L(T; T_0, T_1) = \frac{P(T, T_0)}{P(T, T_1)} - \frac{1}{\tau}
\]
The zero coupon bond $P(T, T_0)$ is an asset, and $P(T, T_1)$ is the numeraire in the $T_1$-forward measure. Thus, FTAP yields that the discounted asset price is a martingale, i.e.,
\[
\mathbb{E}^{T_1} \left[\frac{P(T, T_0)}{P(T, T_1)} \middle| \mathcal{F}_t \right] = \frac{P(t, T_0)}{P(t, T_1)}.
\]

Linearity of the expectation operator yields:
\[
\mathbb{E}^{T_1} \left[L(T; T_0, T_1) \middle| \mathcal{F}_t \right] = \mathbb{E}^{T_1} \left[\frac{P(T, T_0)}{P(T, T_1)} - 1 \middle| \mathcal{F}_t \right] \frac{1}{\tau}.
\]

\[
= \left[\frac{P(t, T_0)}{P(t, T_1)} - 1\right] \frac{1}{\tau}.
\]

\[
= L(t; T_0, T_1).
\]
\end{proof}




\subsection{The Tenor-basis Modelling}


\[
\frac{V(T)}{B(T)} = \mathbb{E}^{\mathbb{Q}} \left[
-\mathbbm{1}_{\{\xi_B > T_0\}} \cdot \frac{1}{B(T_0)} 
+ \mathbbm{1}_{\{\xi_B > T_1\}} \cdot \frac{1 + K \cdot \tau}{B(T_1)}
\;\middle|\; \mathcal{F}_T \right].
\]
\[
\text{(all expectations conditional on } \mathcal{F}_T \text{)}
\]

\[
V(T) = -P(T, T_0) Q(T, T_0) \mathbb{E}^{T_0}[1] 
+ P(T, T_1) Q(T, T_1) \mathbb{E}^{T_1}[1 + K \cdot \tau].
\]

\subsubsection{the hazard rate model}
\begin{itemize}
    \item \textbf{Survival Probability:}
    \[
    S(t) = \mathbb{P}(\tau > t)
    \]
    The probability that the default event has not occurred by time \(t\).

    \item \textbf{Hazard Rate (\(\lambda(t)\)):}
    \[
    \lambda(t) = \lim_{\Delta t \to 0} \frac{\mathbb{P}(t \leq \tau < t + \Delta t \mid \tau > t)}{\Delta t}
    \]
    The instantaneous rate of default at time \(t\), conditional on survival up to \(t\).

    \item \textbf{Relationship Between Hazard Rate and Survival Probability:}
    \[
    S(t) = \exp\left(-\int_0^t \lambda(u) \, du\right)
    \]
    Survival probability decreases exponentially with the cumulative hazard.

    \item \textbf{Applications:}
    \begin{itemize}
        \item \textbf{Credit Default Swap (CDS) Pricing:} Hazard rates are used to model expected cash flows from default protection.
        \item \textbf{Bond Pricing with Default Risk:}
        \[
        P(t) = \int_t^T \exp\left(-\int_t^u (r(s) + \lambda(s)) \, ds\right) c(u) \, du
        \]
        where \(r(t)\) is the risk-free rate, \(\lambda(t)\) is the hazard rate, and \(c(u)\) is the cash flow.
        \item \textbf{Portfolio Credit Risk:} Used to model correlated defaults in portfolios.
    \end{itemize}

    \item \textbf{Advantages:}
    \begin{itemize}
        \item Flexibility to model time-varying hazard rates \(\lambda(t)\).
        \item Simplifies calculations using the exponential relationship with survival probabilities.
        \item Intuitive interpretation of default intensity at any point in time.
    \end{itemize}

    \item \textbf{Challenges and Limitations:}
    \begin{itemize}
        \item Calibration of hazard rates from market or historical data can be difficult.
        \item Assumes independent defaults, which may not hold during crises.
        \item Correlations between multiple entities add complexity.
    \end{itemize}
    
    \item \textbf{Extensions:}
    \begin{itemize}
        \item \textbf{Stochastic Hazard Rates:} Models where \(\lambda(t)\) is itself a stochastic process.
        \item \textbf{Multi-State Hazard Models:} Allows transitions between different credit states (e.g., downgrade, default).
        \item \textbf{Reduced-Form Credit Risk Models:} Focus on the likelihood of credit events rather than explicit asset dynamics.
    \end{itemize}
\end{itemize}


\begin{theorem}
	Deterministic hazard rate assumption preserves the martingale property of the forward Libor rate.
\end{theorem}


\subsection{Vanilla Interest Models}
\subsubsection{Vanilla Interest Rate Option}

\begin{theorem}[Martingale representation theorem]
\end{theorem}

\begin{theorem}[Bachelier Formula]	
\end{theorem}

\begin{theorem}[Black Formula]
\end{theorem}



\section{Introduction To Risk Measures}
\begin{definition}[some properties]\
Given $\mathcal{X}$ the class of financial position, the function $\rho: \mathcal{X} \rightarrow \mathbb{R}$ is a risk measure if it satisfies
\begin{enumerate}
	\item monotonicity: if $X \le Y$, it holds $\rho(Y) \le \rho(X)$. 
	\item translation invariant: $\rho(X + m) = \rho(X) - m$. 
	\item normalisation: $\rho(0) = 0$. 
\end{enumerate}
\end{definition}
\begin{remark}
The cash invariance implies in particular
\begin{align*}
	& \rho(X + \rho(X)) = \rho(X) - \rho(X) = 0.
\end{align*}	
\end{remark}

A risk measure can enjoy many properties, such as convexity (decentralised risk management), positive homogeneous. 

\begin{definition}[acceptance set]
	The acceptance set for the risk  is defined as 
	\begin{align*}
		\mathcal{A}_\rho (X):= \{ X \in \mathcal{X}: \quad \rho(X) \le 0)
	\end{align*}
\end{definition}

\begin{proposition}
	positive homogeneity and convexity imply subadditivity. 
\end{proposition}
Positive homogeneity is strict. In the case where the risk grows nonlinearly w.r.t. position, we only focus on the convexity. In fact, subadditive functions are usually concave, which motivates the following definition:

\begin{definition}[coherent monetary utility functional and coherent risk measure]
A functional $\phi: \mathcal{X} \rightarrow \mathbb{R}$ is called a coherent monetary utility functional if it is monotone, concave, positive homogeneous and cash invariant. The functional 
\begin{align*}
	\rho(X) = - \phi(X),
\end{align*}
is called a coherent risk measure. 
\end{definition}

The following theorem gives a robust representation of the risk measure $\rho$: 
\begin{align*}
	\rho(X) = \sup_{Q \in \mathcal{M}_{1, f}} \bigg( \mathbb{E}^Q(X) - \alpha (Q) \bigg)
\end{align*}

\begin{theorem}
	A convex risk measure $\rho$ on $\mathcal{X}$ is of the form
	\begin{equation}
		\rho(X) = \max_{Q \in \mathcal{M}_{1, f}}  \bigg(\mathbb{E}_Q[-X] - \alpha^{\min}(Q) \bigg),
	\end{equation}
	where
	\begin{equation}
		\alpha^{\min}(Q):= \sup_{X \in \mathcal{A}_p}  \mathbb{E}_Q [-X]  \qquad Q \in \mathcal{M}_{1, f}.
	\end{equation} 
	$\alpha^{\min}$ dominate all the penalty term that admits the representation of $\rho$. 
\end{theorem}
\begin{proof}
\quad\\
\textbf{I.} 
\begin{align*}
	\rho(X) \ge \sup_{Q \in \mathcal{M}_{1, f}} \bigg( \mathbb{E}_Q[-X] - \alpha^{\min} (Q) \bigg)\quad \forall X \in \mathcal{X}
\end{align*}
It holds in particular by the cash invariance
\begin{align*}
	X':= X + \rho(X) \in \mathcal{A}_\rho.
\end{align*}
Taking $X'$ for the $\alpha^{\min}$ function, it holds
\begin{align*}
	\forall Q \in \mathcal{M}_{1, f}: \quad RHS \le  \mathbb{E}_Q[-X] - \mathbb{E}_Q[X] - \rho(X) = \rho(X).
\end{align*}
Hence the statement holds.
\quad\\
\textbf{II.} Next we show the inverse direction, i.e., 
\begin{align*}
	\rho(X) \le \mathbb{E}_{Q_X}[-X] - \alpha^{\min}(Q_X),
\end{align*}
for some $Q_X$ depends on $X$.

\quad\\
\textbf{II.1, simplification} w.o.l.g we prove for $X \in \mathcal{X}$ with $\rho(X) = 0$ and also we assume since $\rho$ is cash invariant. (Why???) And we assume $\rho(0) = 0$. \par 
\quad\par 
Next, define 
\begin{align*}
	\mathcal{B}:= \{ Y \in \mathcal{X}: \rho(Y) < 0\}
\end{align*}
Noted $\mathcal{B}$ is convex and open (why???). Since $\rho(0) = 0$, by Hahn-Banach separation theorem (of open convex subspace) there exist a functional $l$ on $\mathcal{X}$ such that
\begin{align*}
	l(X) \le \inf_{Y \in \mathcal{M}} l(Y):= b. 
\end{align*}
Show the following claim:
\begin{align*}
	& 1). l(Y) \ge 0 \quad \forall Y \in \mathcal{B},\\
	& 2). l(1) > 0. 
\end{align*}
1). By the cash invariance and convexity and $\rho(0) = 0$, it holds
\begin{align*}
	\rho(\lambda Y + (1 - \lambda) \epsilon) & \le \lambda \rho (Y) + ( 1 - \lambda ) \rho(\epsilon) 
\end{align*}
Taking $\epsilon \rightarrow 0$, by continuity it holds for $\lambda > 0$. 
\begin{align*}
	\rho (\lambda Y) \le \lambda \rho(Y) < 0. \Rightarrow \lambda Y \in \mathcal{B}
\end{align*}
By the cash invariance $ 1 + \lambda Y \in \mathcal{B}$, it holds
\begin{align*}
	l( 1 + \lambda Y) = l(1) + \lambda l(Y)
\end{align*}


\quad\\
2) Since $l$ it not trivial on $\mathcal{B}$, there exists $Y \in \mathcal{B}$ with $l(Y) > 0$, normalised $Y$ such that $\|Y\|_{L^\infty} \le 1$(???) and write $l(Y) = l(Y^+) + l(Y^-) \le l(1) $. It holds
\begin{align*}
	& l(1 - Y^+) \ge l(Y^-) > 0, \\
	& l(Y^+) > 0,	
\end{align*}
hence $l(1) = l(1 - Y^+) + l(Y^+) > 0$. By the representation theorem, there exists a content that represent the linear functional $l$, i.e., there exists a $Q \in \mathcal{M}_{1, f}$, such that
\begin{align*}
	l(X) = \mathbb{E}_Q[X] \cdot l(1) \quad \forall X \in \mathcal{X}, 
\end{align*} 
hence
\begin{align*}
	\mathbb{E}_Q  [Y] = \frac{l(Y)}{l(1)} \quad \forall Y \in \mathcal{B}. 
\end{align*}
\end{proof}




\begin{example}[entropic risk measure]
Consider the penalty function defined on the space of measure
\begin{align*}
	\alpha: \mathcal{M} \rightarrow (0, \infty]
\end{align*}
\begin{align*}
	\rho(X):= \sup_{Q \in \mathcal{M}_1} \bigg\{ \mathbb{E}_{\mathbb{Q}} [-X] - \frac{1}{\beta} KL(\mathbb{Q}|\mathbb{P}) \bigg\}.
\end{align*}
\end{example}

\begin{example}[shortfall risk]
	
\end{example}

\subsection{Convex risk measure on $L^\infty$}

\subsection{V$@$R, AV$@$R and shortfall risk.}
The Var Model is one of the most widely used model. It was invented by JP.Morgen, it indicated the quantile of loss, i.e., the loss of the position is at least $X$ with probability at most $\lambda$. It improves the traditional variance-based risk measure in that it calculates the not only the volatility, but also the expected loss under the worst case. On the other hand, it is also criticised as pro-cycle since it's calibrated on the recent data. 
\begin{definition}[value at risk]
	Given the probability space $(\Omega, \mathcal{F}, \mathbb{P})$ and $X$ a random variable on it, we observe the $\lambda$-quantile function
	\begin{align*}
		q_X(\lambda):= \{q \in \Omega:\quad \mathbb{P}[X \le q] \ge \lambda \texttt{ and } \mathbb{P}(X < q) \le \lambda \}. 
	\end{align*}
	By the right-continuity of the cdf we define the value at risk as
	\begin{equation}
		V@R_\lambda(X):= -q^+_X(\lambda). 
	\end{equation}
\end{definition}

\begin{proposition}
	For any $X \in L^\infty$ and each $\lambda \in (0, 1)$, it holds
	\begin{align*}
		V@R_\lambda(X) = \min_\rho \{ \rho(x): \quad \rho \text{ convex,  continuous from above, dominates V@R} \}.  
	\end{align*}
\end{proposition}

\begin{definition}[AV@R]
The average value at risk is defined as
\begin{align*}
	AV@R (X):= \frac{1}{\lambda} \int_0^\lambda V@R_\gamma (X) d\gamma
\end{align*}
\end{definition}


\subsubsection{The Shortfall Risk}
Given $u: \mathbb{R} \rightarrow \mathbb{R}$ the utility function, define the lost function
\begin{align*}
	l(x): \mathbb{R} \rightarrow \mathbb{R}, \quad x \mapsto - u(-x). 
\end{align*}
Define the acceptance set:
\begin{align*}
	\mathcal{A} := \{ X \in \mathcal{X}:  \mathbb{E}[u(x)] \ge y_0 \}.
\end{align*}
We are able to define the convex monetary risk measure $\rho$ (Why??? define as such)
\begin{equation}
	\rho(X):= \inf_{ m \in \mathbb{R}} \{ m |\rho(X + m) \ge y_0 \}
\end{equation}
The following 

\begin{definition}[shortfall risk]
\begin{align*}
	\alpha^{\min}(Q)
\end{align*}
\end{definition}


\chapter{Insurance Mathematics}
\section{Life Insurance}
We start by defining the following notations
\begin{enumerate}
	\item $l_x$: probability of living person at the age of $x$.
	\item ${}_t q_x$: probability of dying of a $x$-years old within $t$ years.
	\item ${}_t p_x$: probability of living of a $x$-years old with $t$ years. 
	\item $v$: technical interest rate
\end{enumerate}

$q^{aa}_x$: died as active

\begin{definition}[life expectancy]
	$\mathbb{E}[T(x)]$ the expected remaining life a $x$-years old person. $T$ is the random variable denoting the remaining life. 
\end{definition}

\begin{definition}[the hazard rate]
	The hazard rate $\mu_{x + t}$ describe the mortality of a given $x$-years old person died at $x + t$:
	\begin{align*}
		\mu_{x + t}:= \lim_{\Delta t \rightarrow 0} \frac{P(t \le T \le T + \Delta t|T \ge t) }{P(T \ge t) }
	\end{align*}
\end{definition}

\begin{lemma}[Chapman-Kolmogorov in mortality]
It holds
\begin{align*}
	{}_tp_x^r = \frac{l_{x + t}}{l_x}
\end{align*}
\end{lemma}
\begin{proof}
\end{proof}

\begin{definition}[commutation value]
The commutation value helps to calculate the present value of the future pension liability. It can be seen as weighted probability of survival.
\begin{align*}
	D_x^s := v^x l_x^s.
\end{align*}
Accordingly we defined the \textbf{discounted remaining life expectancy}
\end{definition}


\quad\\

\subsubsection{PV of the delayed pension and the  }
Notation von Barwert berechnen: 

\begin{enumerate}
	\item $a_x^r$: x-years old, old-age pension, life long
	\item ${}_{z - x} a_x^{a A}$ : x-years old, active to old-age pension at age $z$.   
	\item $ a_{x, z - x\neg}^i$: x-years old, invalid to z. 
	\item 
\end{enumerate}



\begin{lemma}[ewig annuity]
	The present value of the lifelong old-age pension of the age $x$ can be calculated as 
	\begin{align*}
		a_x^r=  \frac{N_x^r}{D_x^r}
	\end{align*}
\end{lemma}
\begin{proof}
It holds
\begin{align*}
	a_x^r &= \sum_{t \ge 0} {}_tp_{x + t}^r \cdot v^t \\
	& = \sum_{t \ge 0} \frac{l_{x + t}}{l_x} \cdot v^t \\
	& = \sum_{ t \ge 0} \frac{l_{x + t} v^{x + t}}{l_x v_x} \\
	& = \frac{N_x^r}{D_x^r}
\end{align*}
\end{proof}

\begin{lemma}[aufgeschobene annuity]
	The invalid pension follows the same formula. However, it is the finite sum since the invalid annuity stops at the age of retirement, i.e., 
	\begin{align*}
		a_x^i = \sum_{t = 0}^{Z - t - 1}
	\end{align*}	
\end{lemma}

\begin{lemma}
	Other calculations
	\begin{enumerate}
		\item ${}_{z - x} a_x^{iA} = ( 1- q_x^a - i_x)(1 - q_{x + 1}^a - i_{x + 1}) \dots (1 - q_{z - 1}^a - i_{z - 1}) \cdot v^{z - x} \cdot a_z^r$
	\end{enumerate}
\end{lemma}


By pricing we established the PV of the insurance contracts at the moment. However, we need to assure the solvency is not jeopardised as the times marches. 



\begin{lemma}[projected Unit Credit Method]
	Denote 
	\begin{enumerate}
		\item $x_0$: age at the beginning of valuation
		\item $k$: Entering data of the benefit
		\item 
	\end{enumerate}
\end{lemma}

\subsection{Valuation Methods}
\subsubsection{Projected Unit Credit Method}
\begin{definition}[Defined (Projected) Benefit Obligation by IFRS (GAAP)]
Present Value of the pension plan at the reference data.
\end{definition}


\begin{proposition}[m/n degressiv method]
	It is calculated as following 
	\begin{enumerate}
		\item Calculate the 
	\end{enumerate}
\end{proposition}



\begin{definition}[covering capital]
\end{definition}


\begin{definition}[embedded value]
\begin{align*}
	EV = PVFP + ADNV
\end{align*}
the adjusted net asset value includes capital gains, tex deduction. 
\end{definition}











\chapter{Population Genetics and Markov Process}
\section{Introduction}
\subsection{The Wright-Fisher random graph}
\begin{definition}[Wright-Fischer random graph]
	Define
	\begin{enumerate}
		\item $V:= \mathbb{Z} \times N$ and $v(g, i)$ with $g \in \mathbb{Z}$, $i \in \{1, \dots, N)$
		\item $\{U_v\}_{v \in V} \overset{d}{\sim} \textbf{U}\{1, \dots, N \}$ are i.i.d random variables
		\item $(g - 1, U_v)$ is the parent of $v$.
		\item The unoriented edge are denoted by $\langle U_v, v \rangle$. The set of the edges is  
			\begin{equation*}
				E:= \{\langle (g_v - 1, U_v), v \rangle | v \in V \}
			\end{equation*}
		\item \textbf{The Wright-Fisher random graph} is represented by 
		 \begin{equation*}
		 	 G := \{V, E\}
		 \end{equation*}
	\end{enumerate}
\end{definition}

\begin{definition}[The natural two-type Wright-Fisher process].
\begin{itemize}
	\item type space of cardinality 2, such as $\{+, - \}$.
	\item $\mu$ a distribution on $\{+ , -\}^N$
	\item The random map $t: V \rightarrow \{+, -\}$ with 
	\begin{enumerate}
		\item The initial distribution $(t(v) | g_{v} = g_0 ) \sim \mu$
		\item $\forall g(v) > g_0: t(v) = t((g(v) - 1, u_v))$
	\end{enumerate}
	\item The \textbf{natural 2-type Wright-Fisher process}
	\begin{equation*}
		X^N(g_0, n) := \frac{1}{N}\sum_{i = 1}^N \textbf{1}\{t(g_0 + n, i) = "-"\}
	\end{equation*}
\end{itemize}
\end{definition}

\begin{remark}.
\begin{itemize}
	\item $X^N(g_0, n)$ is a Markov process because $(U_{v(g, i))})_{g \in \mathbb{Z}}$ are i.i.d. random vector
	\item For the same reason the Markov chain is time-homogeneous and we can write $X^N(g_0, n)$ as
	\begin{equation*}
		X^N(n)
	\end{equation*}
\end{itemize}
\end{remark}

\begin{definition}[Ancester and Common Ancester].\\
	$\tilde{v} \in V$ is an ancester of $v \in V$ if : 
	\begin{itemize}
		\item $g(\tilde{v}) < g(v)$
		\item there exists a genelogical path from $\tilde{v}$ to $v$.
	\end{itemize}
	Define $V_g:= \{v \in V| g(v) > g\}$. We define $\tilde{v} \overset{g}{\sim} v $ if 
	\begin{itemize}
		\item $\tilde{v}, v \in V_g$
		\item $\tilde{v}$, and $v$ have common ancester in generation $g$.
	\end{itemize}
\end{definition}

\begin{theorem}
	$\sim_g$ is an equivalence relation
\end{theorem}
\begin{remark}
$\sim_g$ defines a partition on the set $[N]$.	
\end{remark}

\begin{definition}[The genealogy process]
	Given a WFRG $(V, E)$, we set
	\begin{itemize}
		\item $\mu$ a distribution on $([N], 2^{[N]})$, which is independent from $U_v$. (the sampling)
		\item $S \sim \mu$, $\Pi_0^N := \{\{S\}\}$, $k = | \Pi_0^N|$, $S = [k]$
	\end{itemize}
	The process $(\Pi^N(g_0, n))_{n \in \mathbb{N}_0}$ is the genealogy process of $\Pi^N_0$ given by a partition of $[k]$ induced by $\sim_{g_0 -n}$.
\end{definition}
\begin{remark}
\end{remark}

\begin{definition}[The block-counting process of geneology]
\begin{equation*}
	\forall n \in \mathbb{N}_0: \quad A^N(g_0, n) = |\Pi^N(g_0, n)|
\end{equation*}
\end{definition}

\begin{definition}[Duality]
	The process $(X_i)_{i \in I}$ and $(A_i)_{i \in I}$ is called $H$-dual, if:
	\begin{equation*}
		\forall i \in I, x \in E, a \in F: 
		\mathbb{E}[H(X_i, a) |X_0 = x] = \mathbb{E}[H(x, A_i) | A_0 = a]
	\end{equation*}
	$H$ is called the duality function.
\end{definition}

\begin{theorem}
	\begin{equation*}
		\mathbb{E}[X^N(g_0, n)^a|X^N(g_0, 0) = x] = \mathbb{E}[x^{A^N(g_0, n)}  |A^N(g_0, 0) = a]
	\end{equation*}
	induced by notation in Bayes we write:
	\begin{equation*}
		\mathbb{E}_x[X_n^a] = \mathbb{E}_a[x^{A_n}]
	\end{equation*}
\end{theorem}


\subsection{The scaling limit}
\begin{definition}[the semi-group operator]
Given a state space $E$ and a continuous-time Markov process $X_t$, we define:
\begin{equation*}
	P_t: C^0(E) \rightarrow C^0(E), f \rightarrowtail \mathbb{E}_x[f(X_t)] 
\end{equation*}
as the semi-operator.
\end{definition}

\begin{theorem}[Properties of the semi-group]
	the semi-operator admits the properties of
\begin{itemize}
	\item \textbf{associativity}: $P_t \circ P_s = P_{t + s}$.
	\item \textbf{monoid:} $P_0 = Id_{C_0(E)}$.
\end{itemize}
\end{theorem}
\begin{proof}
	By Chapman-Kolmogrov we have:
	\begin{align*}
		\mathbb{E}_x [\theta_s \circ \theta_t \circ f(X)] 
	   =\mathbb{E}_x[\theta_{s + t} \circ f(X) ]
	\end{align*}
	and 
	\begin{align*}
		\mathbb{E}_x[\theta_s \circ \theta_{-s} \circ f(X)] = \mathbb{E}_x f
	\end{align*}
\end{proof}

\begin{definition}[Generator of the semigroup]
The generator of a semigroup $P_t$ on the state $x$ is defined as
\begin{equation*}
	A(f): D(A) \subset C^0(E) \rightarrow \mathbb{R},\quad f \rightarrowtail \lim_{t \rightarrow 0} \frac{P_t(f)|_x - P_0(f)|_x}{t}
\end{equation*}
or we write
\begin{equation*}
	A(f)|_x = \lim_{t \rightarrow 0} \frac{\mathbb{E}_x[f(X_t) - f(x)]}{t}
\end{equation*}
The domain where $A(f)$ is well-defined is denoted by $D(A)$, it includes the continuous function whose induced semi-group admits a generator.
\end{definition}

\begin{example}[Jump process in finite state space]
Consider $Y_t$ a jump process in finite state space $E$. If the process jumps at exponential time, We have
\begin{align*}
	A(f) &= \lim_{t \rightarrow 0} \frac{\mathbb{E}_y[f(Y_t) - f(y)]}{t}\\
	&= \lim_{t \rightarrow 0} \frac{\sum_{w \in E} \mathbb{P}(N_t = 1)(f(w) - f(y)) r_{yw} }{t}\\
	(1)&= \sum_{w \in E}(f(w) - f(y))r_{wy} \\
	(2)&= \int_E f(w) - f(y) \nu(w)
\end{align*}

(1) follows from that $N_t \sim Poisson(\lambda t)$. Therefore $\frac{\mathbb{P}(N_t = 1)}{t} = \frac{\lambda t }{t} e^{-\lambda t} \overset{ t \rightarrow 0}{\rightarrow}  1 $.\par
(2) we define the measure $\nu(w):= \sum_{w \in E} r_{wy} \delta(y)$
\end{example}

\begin{example}[Deterministic ODE]
Consider $R$ a deterministic process. Furthermore we have the ODE:
\begin{equation*}
	\frac{dR}{dt} = \mu(R)
\end{equation*}
Calculate the generator:
\begin{align*}
	A(f) &= \lim_{t \rightarrow 0} \frac{\mathbb{E}[f(R_t) - f(R_0)]}{t}\\
	&= \lim_{t \rightarrow 0} \frac{f(R_t) - f(R_0)}{t}\\
	&= f'(R_0) \mu(R_0)
\end{align*}
\end{example}

\begin{example}[Random Walk]
Given a Poisson process $N_t$, we can write the generator:
\begin{equation*}
	A(f)|_N = N (f(n+1) - f(n)) 
\end{equation*}
A random walk admits the generator:
\begin{equation*}
	A(f)|_x = r(x)(f(x+1) - f(x)) + q(x)(f(x-1) - f(x))
\end{equation*}
If the random walk is \textbf{symmetric} indexed by the Poisson process, then the jumping time of the random walk is controlled by the Poisson process. The generator of the process $(S_{N_t})_t$ could be written as
\begin{align*}
	A(f)|_{S_{N_t}} = \frac{N}{2} (f(x+1) - f(x)) + \frac{N}{2} (f(x-1) - f(x))
\end{align*} 
By rescaling the step size and applying the Taylor theorem, given $f \in C^3(\Omega)$, the process $\frac{S_{N_t}}{\sqrt{N}}$ admits the generator:
\begin{align*}
	A(f) &= \frac{N}{2}(\frac{f''(x)}{N} + o(N^{-\frac{3}{2}}))\\
	&= \frac{f''(x)}{2} + o(N^{-\frac{1}{2}})
\end{align*}
Furthermore, the speed of the up-down movement could be controlled by the position of $x$, therefore we have:
\begin{align*}
	A(f) = \frac{f''(x) \sigma^2(x)}{2} + o(N^{\frac{1}{2}}) \quad
	 \overset{N \rightarrow \infty}{\longrightarrow} \quad
	 \frac{f''(x) \sigma^2(x)}{2}
\end{align*} 
The greater the $\sigma^2(x)$, the faster the Brownian motion moves.
\end{example}

\begin{definition}[Convergence of Operator]
	A sequence of stochastic process $(S_t^N)_{t \in T}$ converges in operator if 
	\begin{equation*}
		\forall f \in D(A): A(f)_{S_N} \overset{L^1}{\rightarrow} A(f)_S
	\end{equation*}
\end{definition}
\begin{remark}
	
\end{remark}


\begin{theorem}
\end{theorem}


\subsection{Wright-Fisher diffusion}
\begin{definition}[Diffusion]
\end{definition}

\begin{theorem}[Ito Formula]
\end{theorem}


\end{document}



































